\documentclass[11pt,letterpaper]{article}

\usepackage[spanish,es-noshorthands]{babel}
\usepackage{fontspec}
\setmainfont{Source Serif Pro}
\usepackage{microtype}
\usepackage{geometry}
\usepackage{enumitem}
\usepackage{xcolor}
\usepackage{csquotes}
\usepackage{float}
\usepackage{tikz}
\usetikzlibrary{arrows.meta,positioning,fit,backgrounds}
\usepackage{xurl}
\usepackage[hidelinks]{hyperref}
\usepackage[backend=biber,style=apa,sorting=nyt,maxcitenames=2]{biblatex}
\DeclareLanguageMapping{spanish}{spanish-apa}

\geometry{margin=2.3cm}
\hbadness=10000
\setlength{\emergencystretch}{1em}
\setlength{\parindent}{0pt}
\setlength{\parskip}{0.5em}
\setlist{nosep}
\addbibresource{referencias.bib}
\AtBeginBibliography{\sloppy\small}
\setlength{\bibitemsep}{0.3em}

% Colores: verde para los avances de la ciencia, terracota para los hitos de
% gobernanza, azul oscuro para el eje del tiempo y gris para el texto
% secundario.
\definecolor{cCiencia}{RGB}{36,122,82}
\definecolor{cGob}{RGB}{186,86,30}
\definecolor{cEje}{RGB}{26,58,102}
\definecolor{cGris}{RGB}{120,116,110}

\tikzset{
  % los rótulos de la figura no se dividen con guion
  every node/.append style={execute at begin node={\hyphenpenalty=10000}},
  hito/.style={draw=#1, fill=#1!9, thick, rounded corners=2pt,
    text width=1.74cm, align=center, inner sep=2.5pt, font=\scriptsize},
  anio/.style={font=\footnotesize\bfseries, text=cEje},
  punto/.style={circle, fill=#1, inner sep=0pt, minimum size=5pt},
  carril/.style={font=\footnotesize\bfseries, text=#1, anchor=west},
}

\begin{document}

% --- Portada -----------------------------------------------------------------
\begin{titlepage}
\centering
\vspace*{3.2cm}
{\large\scshape Gestión del Desarrollo Tecnológico:\\
Estrategias y Análisis de Portfolio\par}
\vspace{0.5cm}
{\normalsize Módulo 5. Foro 5.1: Bioética\par}
\vspace{2.6cm}
{\color{cEje}\rule{0.55\textwidth}{0.8pt}\par}
\vspace{0.7cm}
{\LARGE\bfseries Entender la vida\\[0.2em]
antes de rediseñarla\par}
\vspace{0.5cm}
{\large Por qué la biotecnología necesita más herramientas\\
y estándares comunes, no prohibiciones\par}
\vspace{0.7cm}
{\color{cEje}\rule{0.55\textwidth}{0.8pt}\par}
\vfill
{\large Carlos Luis Rojas Aragonés\par}
\vspace{0.3cm}
{\normalsize Chief Technology Officer Program\par}
\vspace{0.3cm}
{\normalsize 5 de octubre de 2026\par}
\vspace*{1.5cm}
\end{titlepage}

% --- Cuerpo ------------------------------------------------------------------
\section*{1. La vida sigue siendo un misterio}

Cada organismo vivo es, por sí mismo, extraordinario. En \emph{The Body},
\textcite{bryson2019} recuerda que en la escuela le enseñaron que los
elementos químicos de un cuerpo humano costaban unos pocos dólares; la Royal
Society of Chemistry calculó en 2013 que en realidad rondan los 151\,000
dólares. Aun así, ni los científicos más brillantes del planeta podrían
ensamblar con ellos una persona. Ni siquiera una célula: el genoma sintético
de \textcite{gibson2010} tuvo que instalarse en una bacteria que ya existía.

\section*{2. Usar todas las herramientas a nuestro alcance}

Por eso creo que debemos aprovechar toda herramienta que nos ayude a entender
los organismos. AlphaFold resolvió un problema de cincuenta años, predecir la
forma de las proteínas, y les valió a Hassabis y Jumper el Nobel de Química
2024, compartido con David Baker por diseñar proteínas nuevas
\parencite{nobel2024}. Los concursos abiertos aceleran este avance: en el de
Adaptyv sobre EGFR, 130 equipos enviaron 1\,131 diseños y el ganador se unió a
su objetivo ocho veces más fuerte que el anticuerpo comercial cetuximab
\parencite{cotet2025}. Hoy Anthropic y Adaptyv financian la validación de más
de 5\,000 diseños contra cinco enfermedades y publicarán todos los resultados,
incluidos los negativos \parencite{proteinbase2026}.

\section*{3. Límites sí, pero como estándar mundial}

No propongo eliminar las regulaciones, sino estandarizarlas como práctica
mundial, en la línea del marco de la \textcite{oms2021}. El caso de He Jiankui,
que en 2018 editó embriones humanos \parencite{cyranoski2018}, muestra el
riesgo de que alguien actúe por su cuenta. Aun con ese riesgo, el balance de
CRISPR, Nobel 2020 para Charpentier y Doudna \parencite{nobel2020}, es
claramente positivo:

\begin{itemize}[leftmargin=1.2em, itemsep=0.25em]
  \item \textbf{Casgevy (2023).} 29 de 31 pacientes con anemia falciforme
    quedaron sin crisis de dolor \parencite{fda2023}.
  \item \textbf{Bebé KJ (2025).} Primera edición genética diseñada para un solo
    paciente, en seis meses \parencite{musunuru2025}.
\end{itemize}

Lo resume la Figura~\ref{fig:linea}.

\begin{figure}[H]
\centering
\begin{tikzpicture}[x=1.88cm]
  % Eje del tiempo
  \draw[very thick, cEje, -{Stealth[length=2.8mm]}] (-0.45,0) -- (8.5,0);

  % Hitos de la ciencia, sobre el eje
  \foreach \x/\a/\t in {
      0/2010/{Genoma sintético en una célula existente},
      1/2012/{Se publica CRISPR programable},
      3/2020/{AlphaFold 2 y Nobel a CRISPR},
      5/2023/{Casgevy, primera terapia CRISPR},
      6/2024/{Nobel a AlphaFold y al diseño de proteínas},
      7/2025/{Bebé KJ: edición a la medida},
      8/2026/{Anthropic × Adaptyv: 5\,000 diseños abiertos}} {
    \node[punto=cCiencia] at (\x,0) {};
    \draw[cCiencia, thick] (\x,0) -- (\x,0.42);
    \node[hito=cCiencia, anchor=south] at (\x,0.42) {\t};
    \node[anio, anchor=north] at (\x,-0.12) {\a};
  }

  % Hitos de la gobernanza, bajo el eje
  \foreach \x/\a/\t in {
      2/2018/{He Jiankui edita embriones humanos},
      4/2021/{Marco de gobernanza de la OMS}} {
    \node[punto=cGob] at (\x,0) {};
    \node[anio, anchor=south] at (\x,0.12) {\a};
    \draw[cGob, thick] (\x,0) -- (\x,-0.42);
    \node[hito=cGob, anchor=north] at (\x,-0.42) {\t};
  }

  % Carriles
  \node[carril=cCiencia] at (-0.45,2.35) {Ciencia y tecnología};
  \node[carril=cGob]     at (5.3,-1.15) {Gobernanza y ética};
\end{tikzpicture}
\caption{La ciencia avanza y la gobernanza la acompaña. Sobre el eje, los
avances que nos ayudan a entender y corregir la vida
\parencite{gibson2010,jinek2012,jumper2021,fda2023,nobel2024,musunuru2025,proteinbase2026};
bajo el eje, los hitos que muestran por qué hacen falta reglas comunes
\parencite{cyranoski2018,oms2021}.}
\label{fig:linea}
\end{figure}

\printbibliography[title={Referencias}]

\end{document}
